\section{Exact likelihood, mixed-rate LAN, and efficiency}\label{sec:lan}

Let $\mathcal E_{N,n}$ be the experiment generated by the transitions in \eqref{eq:transition}. Conditional on the initial values, the log-likelihood is
\[
 \ell_{N,n}(\theta,\sigma)
 =-\frac12\sum_{k=1}^N\sum_{j=0}^{n-1}
 \left\{
 \log\eta_{k,n}^2
 +\frac{(Y_{k,j+1}-\rho_{k,n}Y_{k,j})^2}{\eta_{k,n}^2}
 \right\}+C.
\]
The scale parameter can be profiled explicitly for every $\theta$, leaving a one-dimensional compact optimisation.

Set
\[
 \Lambda_N=\sum_{k=1}^N\lambda_k,
 \qquad
 A_{N,n}=\diag(\sqrt{\Lambda_N},\sqrt{Nn}).
\]

\begin{assumption}[Joint spectral and temporal regime]\label{ass:regime}
As $N,n\to\infty$,
\[
 \lambda_N\Delta_n\to0,
 \qquad
 \frac{\lambda_N}{\Lambda_N}\to0,
\]
and the deterministic initial energy satisfies
\[
 \sum_{k\le N}\frac{\lambda_ku_{0,k}^2}{q_k^2}=o(\Lambda_N).
\]
\end{assumption}

\begin{theorem}[Mixed-rate LAN]\label{thm:lan}
Under \Cref{ass:core,ass:regime}, for every bounded $h\in\R^2$,
\[
 \ell_{N,n}(\vartheta_0+A_{N,n}^{-1}h)-\ell_{N,n}(\vartheta_0)
 =h^\top\Delta_{N,n}-\frac12h^\top J(\vartheta_0)h+o_{\Pp}(1),
\]
uniformly on compact $h$-sets, where
\[
 \Delta_{N,n}\Rightarrow N(0,J(\vartheta_0)),
 \qquad
 J(\vartheta_0)=
 \begin{pmatrix}
 T/(2\theta_0)&0\\
 0&2/\sigma_0^2
 \end{pmatrix}.
\]
The exact conditional MLE is consistent, asymptotically linear, and efficient:
\[
 \begin{pmatrix}
 \sqrt{\Lambda_N}(\widehat\theta-\theta_0)\\
 \sqrt{Nn}(\widehat\sigma-\sigma_0)
 \end{pmatrix}
 \Rightarrow
 N\left(
 0,
 \begin{pmatrix}
 2\theta_0/T&0\\
 0&\sigma_0^2/2
 \end{pmatrix}
 \right).
\]
Polynomial likelihood-ratio bounds imply convergence of polynomial moments.
\end{theorem}

\inwords{Spatial replication supplies the drift information that fixed-horizon high-frequency sampling of one diffusion path cannot. The scale accumulates information from every transition and is learned faster.}

For the Dirichlet heat equation on $(0,1)$,
\[
 \lambda_k=\pi^2k^2,
 \qquad
 \Lambda_N\sim\frac{\pi^2}{3}N^3.
\]
The temporal condition is $N^2/n\to0$, and
\[
 N^{3/2}(\widehat\theta-\theta_0)
 \Rightarrow
 N\left(0,\frac{6\theta_0}{\pi^2T}\right).
\]

\begin{remark}[Where the drift information comes from]\label{rem:info-arith}
The rate $\sqrt{\Lambda_N}$ has a two-line derivation. One transition of mode $k$ carries drift information $(\partial_\theta\rho_k)^2\E Y_{k,j}^2/\eta_k^2\approx(\lambda_k\Delta)^2\cdot\frac{\sigma^2q_k^2}{2\theta\lambda_k}\cdot\frac1{\sigma^2q_k^2\Delta}=\frac{\lambda_k\Delta}{2\theta}$ at stationarity; over $n$ transitions this is $\lambda_kT/(2\theta)$ per mode --- \emph{finite}, because a fixed-horizon path of one mode cannot identify $\theta$ --- and summing across observed modes gives $\Lambda_N\,T/(2\theta)$. Spatial replication, not time, drives drift information; the scale $\sigma$, in contrast, collects $2/\sigma^2$ from every one of the $Nn$ standardised innovations. This is the entire content of the limit $J(\vartheta_0)$ in \Cref{thm:lan}; the proof in part S3 controls the remainder of this arithmetic uniformly.
\end{remark}

\begin{remark}[Rate separation within the joint regime]\label{rem:cross}
The normalised drift--scale cross-information vanishes at rate
$\sqrt{\Lambda_N/(Nn)}\le\sqrt{\lambda_N\Delta_n/T}$, since
$\Lambda_N\le N\lambda_N$ and $n\Delta_n=T$; asymptotic orthogonality of
the two coordinates is therefore a consequence of resolving the fastest
observed relaxation time, not an additional assumption. Complete proofs of
\Cref{thm:lan} --- the score CLT, the uniform quadratic remainder,
consistency via exact scale profiling, and a polynomial likelihood-ratio
bound yielding moment convergence --- are given in \ifdefined\MERGEDDOC Part~S3 (\Cref{mpart:wp3})\else part S3 of the technical supplement\fi.
\end{remark}

\begin{figure}[t]
\centering
\resizebox{.9\textwidth}{!}{\begin{tikzpicture}[font=\small,
  ax/.style={-{Stealth[length=2.6mm]}, thick},
  lbl/.style={font=\footnotesize\itshape, text=black!70}]

% mapping: X = xa*(log10 N - 0.4), Y = ya*(log10 n - 2.5); ranges logN in [0.4,1.8], logn in [2.5,5.3]
\pgfmathsetmacro{\xa}{5.5}\pgfmathsetmacro{\ya}{2.4}
\pgfmathsetmacro{\Xmax}{\xa*1.4}   % logN=1.8
\pgfmathsetmacro{\Ymax}{\ya*2.8}   % logn=5.3

% boundary line: log n = 0.903 + 2 log N  (lambda_N Delta_n = const)
% at logN=0.4 -> logn=1.703 (below plot, clip); crosses logn=2.5 at logN=0.7985
\pgfmathsetmacro{\xcross}{\xa*(0.7985-0.4)}
\pgfmathsetmacro{\ybndR}{\ya*(0.903+2*1.8-2.5)}  % at logN=1.8 -> logn=4.503

% shading
\fill[blue!8] (0,0) -- (\xcross,0) -- (\Xmax,\ybndR) -- (\Xmax,\Ymax) -- (0,\Ymax) -- cycle;
\fill[red!8]  (\xcross,0) -- (\Xmax,0) -- (\Xmax,\ybndR) -- cycle;

% axes
\draw[ax] (0,0) -- (\Xmax+0.55,0);
\node[anchor=north] at (\Xmax*0.5,-0.55) {$\log_{10}N$ \ (Fourier modes)};
\draw[ax] (0,0) -- (0,\Ymax+0.45);
\node[anchor=south, rotate=90] at (-0.85,\Ymax*0.5) {$\log_{10}n$ \ (time steps)};

% ticks
\foreach \N/\lx in {5/0.699, 10/1.0, 20/1.301, 40/1.602}{
  \draw ({\xa*(\lx-0.4)},0) -- ++(0,-0.09) node[below, font=\scriptsize]{\N};}
\foreach \n/\ly in {1000/3.0, 8000/3.903, 64000/4.806}{
  \draw (0,{\ya*(\ly-2.5)}) -- ++(-0.09,0) node[left, font=\scriptsize]{\n};}

% boundary line
\draw[dashed, thick, black!60] (\xcross,0) -- (\Xmax,\ybndR)
  node[pos=.62, below, sloped, lbl]{$\lambda_N\Delta_n\asymp1$ \ ($n\asymp N^2$)};

% design line n = 80 N^2: log n = 1.903 + 2 log N; from logN=0.5 to logN=1.55
\draw[thick, blue!60!black]
  ({\xa*(0.5-0.4)},{\ya*(1.903+2*0.5-2.5)}) -- ({\xa*(1.45-0.4)},{\ya*(1.903+2*1.45-2.5)})
  node[pos=.45, above, sloped, font=\footnotesize, blue!60!black]{$n=80N^2$};

% design points, labels below-right to clear the line
\foreach \N/\n/\lx/\ly in {5/2000/0.699/3.301, 10/8000/1.0/3.903, 20/32000/1.301/4.505}{
  \fill[blue!60!black] ({\xa*(\lx-0.4)},{\ya*(\ly-2.5)}) circle (2.4pt);
  \node[below right=0pt and 1pt, font=\scriptsize, blue!60!black]
    at ({\xa*(\lx-0.4)},{\ya*(\ly-2.5)}) {$(\N,\n)$};}

% regime labels
\node[align=left, font=\footnotesize, blue!50!black, anchor=north west] at (0.25,\Ymax-0.15)
  {\textbf{theorem-aligned regime:} $\lambda_N\Delta_n\to0$ \ (LAN holds)\\[2pt]
   $\sqrt{\Lambda_N}\,(\widehat\theta-\theta_0)\Rightarrow N\!\big(0,\tfrac{2\theta_0}{T}\big)$\qquad
   $\sqrt{Nn}\,(\widehat\sigma-\sigma_0)\Rightarrow N\!\big(0,\tfrac{\sigma_0^2}{2}\big)$};
\node[align=center, font=\footnotesize, red!55!black, anchor=north] at (\Xmax-1.35,1.95)
  {$\lambda_N\Delta_n\not\to0$:\\ fast modes under-resolved,\\ no LAN guarantee};

% side annotation: rate anisotropy (outside plot, right)
\node[align=left, font=\footnotesize, draw=black!30, rounded corners=2pt, fill=white,
      inner sep=5pt, anchor=west] at (\Xmax+0.75,\Ymax*0.62)
  {mixed rates (heat eq., $d=1$):\\[2pt]
   drift\ \ $\sqrt{\Lambda_N}\asymp N^{3/2}$\ \ (spatial replication)\\[1pt]
   scale\ \ $\sqrt{Nn}$\ \ (every transition)\\[3pt]
   cross-term control:\\[1pt]
   $\dfrac{\Lambda_N}{Nn}\le\dfrac{\lambda_N\Delta_n}{T}\to0$};
\end{tikzpicture}
}
\caption{Theorem-aligned designs for the Dirichlet heat equation. The publication designs sit on $n=80N^2$, inside the regime $\lambda_N\Delta_n\to0$ where \Cref{thm:lan} applies; drift and scale are learned at the anisotropic rates shown.}
\label{fig:regimes}
\end{figure}

Continuous spectral observation is not used as the joint regular experiment because quadratic variation reveals $\sigma$ exactly on a continuously observed path. The discrete experiment retains a regular two-parameter likelihood.

\begin{remark}[Rate comparison with continuous observation]\label{rem:hr-compare}
Under continuous observation of the first $N$ modes on $[0,T]$, the drift MLE attains the same spectral rate $\sqrt{\Lambda_N}$ \citep{HuebnerRozovskii1995}; see \citet{Lototsky2009,Cialenco2018}. The contribution of discreteness here is not a different drift rate but a \emph{regular joint} experiment: $\sigma$ becomes a genuine parameter learned at $\sqrt{Nn}$, and the LAN structure with diagonal normalisation transfers to every post-intervention functional in \Cref{sec:uniform}.
\end{remark}
